%% This document is created by 
%%  Dr. Putu Harry Gunawan
%% Template untuk Proposal TA 1 dan TA
%% Template ini digunakan untuk penulisan proposal TA 1 atau TA Fakultas Informatika, Telkom University.

\documentclass[a4paper,12pt,oneside]{book}
\usepackage[utf8]{inputenc}
\usepackage{sectsty}
\usepackage{graphicx}
\usepackage{epstopdf}
\usepackage{algorithm}
\usepackage{algpseudocode}
\usepackage{array}
\usepackage[table]{xcolor}
\usepackage{anysize}
\usepackage{amsmath}
\usepackage{amssymb}
\usepackage[bahasa]{babel}
\usepackage{indentfirst} %Spasi untuk paragraf pertama
\usepackage{geometry}
\usepackage{multirow}% http://ctan.org/pkg/multirow
\usepackage{hhline}% http://ctan.org/pkg/hhline
\marginsize{4cm}{3cm}{3cm}{3cm} %{left}{right}{top}{bottom}
\usepackage[compact]{titlesec} 
\usepackage{etoolbox}

\makeatletter
\patchcmd{\ttlh@hang}{\parindent\z@}{\parindent\z@\leavevmode}{}{}
\patchcmd{\ttlh@hang}{\noindent}{}{}{}
\makeatother

\chapterfont{\centering}
\newcommand{\bigsize}{\fontsize{16pt}{14pt}\selectfont}
\chapterfont{\centering\bigsize\bfseries}
\sectionfont{\large\bfseries}
\usepackage{tikz}
\usetikzlibrary{shapes.geometric, arrows}
%\renewcommand{\chaptertitle}{BAB}
\renewcommand{\thechapter}{\Roman{chapter}}
\renewcommand\thesection{\arabic{chapter}.\arabic{section}}
\renewcommand\thesubsection{\thesection.\arabic{subsection}}
\renewcommand{\theequation}{\arabic{chapter}.\arabic{equation}}
\renewcommand{\thefigure}{\arabic{chapter}.\arabic{figure}}
\renewcommand{\thetable}{\arabic{chapter}.\arabic{table}}

\renewcommand\bibname{Daftar Pustaka}
\addto{\captionsbahasa}{\renewcommand{\bibname}{Daftar Pustaka}}
\usepackage{fancyhdr}
\pagestyle{fancy}
\lhead{}
\chead{}
\rhead{}
\lfoot{}
\cfoot{\thepage}
\rfoot{}
\renewcommand{\headrulewidth}{0pt}

\makeatletter

%%%%%%%%%%%%%%%%%%%%%%%%%%%%%%%%%%%%%%%%%%%%%%%%%%%%%%%%%%%%
%
%  Berikut adalah data-data yang wajib diisi oleh mahasiswa
%
%%%%%%%%%%%%%%%%%%%%%%%%%%%%%%%%%%%%%%%%%%%%%%%%%%%%%%%%%%%%

\title{Keberadaan Solusi Tunggal Persamaan Gelombang Akustik pada Senar Gitar dan Ukulele}\let\Title\@title   %Judul dalam bahasa Indonesia

\newcommand{\EngTitle}{The existence of unique solution of Wave Acoustic Equation on Guitar and Ukulele String}  %Judul dalam bahasa Inggris

%\author{Bang Roma Irama}  \let\Author\@author  %Nama mhs
%\newcommand{\NIM}{30111002}
\newcommand{\Prodi}{Teknik Informatikai}
\newcommand{\KK}{Modelling Computational Experiment} %UNTUK TA
\newcommand{\Gelar}{Sains Komputasi} % UNTUK TA
\date{2016}           \let\Date\@date %Maskkan hanya tahun saja
\newcommand{\Tanggal}{22} % Tanggal Pengesahan
\newcommand{\Bulan}{Agustus} % Bulan Pengesahan
\newcommand{\PembimbingSatu}{Dr. Calon Pembimbing 1, M.Kom}
\newcommand{\NIPPembimbingSatu}{123456}
\newcommand{\PembimbingDua}{Dr. Calon Pembimbing 2, M.Kom}
\newcommand{\NIPPembimbingDua}{123456}
\newcommand{\Kaprodi}{Dr. Ketua Program Studi}
\newcommand{\NIPKaprodi}{95650581-1}
\newif\iflogTA
%\logTAtrue   %%%%%% WARNING kode ini diaktifkan untuk format TUGAS AKHIR
\makeatother
\linespread{1}


\begin{document}
\pagenumbering{roman} 
%%\maketitle
\begin{titlepage}
\thispagestyle{empty}
%\vspace*{0.7cm}
{\centering
\large
{\bigsize\bf \Title}\\
\vspace{ 2cm}
\rm
\iflogTA
\textbf{Tugas Besar}\\
\vspace{0.5 cm}
\textbf{Kelompok Keahlian: \KK}\\
\else
\textbf{Laporan Tugas Besar Penambangan Data}\\
\vspace{0.5 cm}
\textbf{SIDE - 02}\\
\fi
\vspace{0.5 cm}
\textbf{1103134344      Zavli Juwantara}


\textbf{1103130277      Bayu Rizky Ramadhan}


\textbf{1301140000      Imam Rahman}
\vspace{1.5 cm}

\begin{figure}[h]
{\centering {\includegraphics[scale=0.17]{Tel-U-Logo}}\par}
\end{figure}

\vspace{2 cm}
{\bigsize\textbf{Program Studi Sarjana \Prodi}\\
\vspace{0.5 cm}
\textbf{Fakultas Informatika}\\
\vspace{0.5 cm}
\textbf{Universitas Telkom}\\
\vspace{0.5 cm}
\textbf{Bandung}\\
\vspace{0.5 cm}
\textbf{\Date}\\}
}

\pagebreak
\thispagestyle{empty}
%{\centering
\iflogTA
\textbf{\large Lembar Pengesahan}\\  %UNTUK TA
\else
\textbf{\large Lembar Persetujuan}\\
\fi
\vspace{0.5cm}
\textbf{\Title}\\
\vspace{0.5cm}
\textbf{\textit{\EngTitle}}\\
\vspace{0.5cm}
%\textbf{\Author}\\
%\textbf{NIM: \NIM}\\
\vspace{1cm}

\iflogTA
{ Tugas Akhir ini diterima dan disahkan untuk memenuhi sebagian dari syarat untuk memperoleh gelar sarjana \Gelar\\ Program Studi Sarjana \Prodi\\ Fakultas Informatika Universitas Telkom}\\  %% UNTUK TA
\else
{ Proposal ini diajukan sebagai usulan pembuatan tugas akhir pada\\ Program Studi Sarjana \Prodi\\ Fakultas Informatika Universitas Telkom}\\
\fi
\vspace{0.5cm}

{Bandung, \Tanggal\quad \Bulan \quad \Date}\\
{Menyetujui}\\

\vspace{0.5cm}
\iflogTA
\begin{center}
\begin{tabular}{  m{8cm}  m{8cm} }
Pembimbing 1 & Pembimbing 2
\end{tabular}
\end{center}
\else
\begin{center}
\begin{tabular}{  m{8cm}  m{8cm} }
Calon Pembimbing 1 & Calon Pembimbing 2
\end{tabular}
\end{center}
\fi
\begin{center}
\vspace{2cm}
\begin{tabular}{  m{8cm}  m{8cm} }
\underline{\PembimbingSatu} & \underline{\PembimbingDua} \\
NIP: \NIPPembimbingSatu & NIP: \NIPPembimbingDua
\end{tabular}
\end{center}
\vspace{0.5cm}
\iflogTA
Mengesahkan,\\   %% UNTUK TA
Kepala Program Studi \Prodi\\ %% UNTUK TA
\vspace{2.5cm}   %% UNTUK TA
\underline{\Kaprodi}\\ NIP: \NIPKaprodi\\  %% UNTUK TA
\fi
}

\pagebreak
\end{titlepage}
\addcontentsline{toc}{chapter}{Abstrak}
%\chapter*{Abstrak}

 Abstrak Bahasa Indonesia
  
\vspace{0.5 cm}
\begin{flushleft}
{\textbf{Kata Kunci:} Shallow, water, equations.}
\end{flushleft}
\iflogTA
\pagebreak
\addcontentsline{toc}{chapter}{Abstract}
\chapter*{Abstract}

  My abstract here

\vspace{0.5 cm}
\begin{flushleft}
{\textbf{Keywords:} Shallow, water, equations.}
\end{flushleft}
\pagebreak
\addcontentsline{toc}{chapter}{Lembar Persembahan}
\chapter*{Lembar Persembahan}

  Sata persembahkan
  \begin{enumerate}
      \item Kepada Tuhan Yang Maha Esa
      \item Kepada Kedua Orangtua saya
      \item Kepada Rekan kerja dll
  \end{enumerate}
\pagebreak
\addcontentsline{toc}{chapter}{Kata Pengantar}
\chapter*{Kata Pengantar}

Puja dan puji syukur saya panjatkan...
\pagebreak
\fi
\cleardoublepage
\addcontentsline{toc}{chapter}{Daftar Isi}
\tableofcontents
\iflogTA
\newpage
\cleardoublepage
\addcontentsline{toc}{chapter}{Daftar Gambar}
\listoffigures
\newpage
\cleardoublepage
\addcontentsline{toc}{chapter}{Daftar Tabel}
\listoftables
%\pagebreak
\fi
%
\cleardoublepage
\pagenumbering{arabic}
\chapter{Pendahuluan}

\section{Deskripsi Data}

Dataset yang digunakan pada Tugas Besar Mata Kuliah Penambangan Data ini adalah dataset dari KDD CUP 1999. Pada dataset tersebut peserta KDD CUP diberikan tugas untuk membangun sebuah \emph{Network Intrusion Detector}, model prediktif yang mampu dalam membedakan antara {\emph{"bad" connections}} yang bisa disebut sebagai gangguan ataupun serangan dan juga {\emph{"good" normal connections}}. \emph{Database} ini berisi satu set  data standar yang akan harus diaudit dengan cakupan berbagai gangguan yang disimulasikan dalam lingkungan jaringan militer.

%
\chapter{Kajian Pustaka}

\section{Persamaan Air Dangkal}
Berikut diberikan persamaan pengatur dari persamaan gelombang pada gitar
\begin{equation}\label{nama-rumus}
    \int_0^1 \frac{f(x)}{g(x)}\ {\rm dx}=\sin x
\end{equation}

Rumus (\ref{nama-rumus}) merupakan contoh persamaan matematika. persamaan matematika diatas diberi nama \textbackslash label\{nama-rumus\}.

\begin{figure}[h!]
    \centering
    \includegraphics[scale=0.3]{Tel-U-Logo.png}
    \caption{Caption}
    \label{fig:my_label}
\end{figure}
\subsection{Cara memanggil pustaka}
Contoh pustaka prosiding \cite{doyen2014explicit}, jurnal \cite{gunawan2015hydrostatic} dan buku \cite{toro2013riemann}. Atau dapat juga mengguanakan dua pustaka atau lebih dalam \cite{gunawan2015hydrostatic,toro2013riemann}.
%
\chapter{Metodologi dan Desain Sistem}
\section{Flowchart sistem}

%Definisikan bentuk dan warna
\tikzstyle{startstop} = [rectangle, rounded corners, minimum width=3cm, minimum height=1cm,text centered, draw=black, fill=red!30]
\tikzstyle{io} = [trapezium, trapezium left angle=70, trapezium right angle=110, minimum width=3cm, minimum height=1cm, text centered, draw=black, fill=blue!30]
\tikzstyle{process} = [rectangle, minimum width=3cm, minimum height=1cm, text centered, minimum width=3cm, draw=black, fill=orange!30]
\tikzstyle{decision} = [diamond, minimum width=3cm, minimum height=1cm, text centered, draw=black, fill=green!30]
\tikzstyle{arrow} = [thick,->,>=stealth]
\tikzstyle{cloud} = [draw, ellipse,fill=red!20,
    minimum height=2em]

\begin{figure}[h!]
    \centering
    %Mulai menggambar Flowchart
\begin{tikzpicture}[node distance=2cm]
\node (start) [cloud] {Start};
\node (in1) [io, below of=start] {Input};
\node (pro1) [process, below of=in1] {Process 1};
\node (dec1) [decision, below of=pro1] {Decision 1};
\node (pro2a) [process, below of=dec1, yshift=-0.5cm] {Process 2a};
\node (pro2b) [process, right of=dec1, xshift=2cm] {Process 2b};
\node (out1) [io, below of=pro2a] {Output};
\node (stop) [cloud, below of=out1] {Stop};
\draw [arrow] (start) -- (in1);
\draw [arrow] (in1) -- (pro1);
\draw [arrow] (pro1) -- (dec1);
\draw [arrow] (dec1) -- (pro2a);
\draw [arrow] (dec1) -- (pro2b);
\draw [arrow] (pro2b) |- (pro1);
\draw [arrow] (pro2a) -- (out1);
\draw [arrow] (out1) -- (stop);
\end{tikzpicture}
    \caption{Caption flowchart}
    \label{figflow}
\end{figure}

\section{Algoritma}
 Atau dalam bentuk algoritma seperti contoh pada Algoritma \ref{Algo:FVDM} berikut ini:
 

\begin{algorithm}
 \begin{algorithmic}[1]
    \Procedure{FVDM}{$Tfinal, \Delta t$}
    \State \text{Start}
	\State \textbf{For }$n=1:N$\textbf{ do} \Comment{Pemberian nilai awal}
	\State \hspace{0.5cm} Input nilai $x[n]$
	\State \hspace{0.5cm} Input nilai $v[n]$
	\State \textbf{EndFor}
	\State \text{time=0}
	\While{$time < Tfinal$}
	\State \hspace{0.5cm} $time=time +\Delta t$
	\State \hspace{0.5cm} Hitung jarak bamper menggunakan rumus  untuk $n=2,\cdots,N$ 
	\State \hspace{0.5cm} \textbf{If}( $S(n) \leq 0 m)$ \textbf{then return End If}.
	\State \hspace{0.5cm} Tentukan $\lambda$ menggunakan.
	\State \hspace{0.5cm} Hitung kecepatan optimal $v_o(t)$ menggunakan.
	\State \hspace{0.5cm} Hitung percepatan $a_n(time)$ menggunakan .
	\State \hspace{0.5cm} Hitung kecepatan baru dengan $v_n(time)=v_n(time-\Delta t) + a_n(time) \Delta t$.
	\State \hspace{0.5cm} Hitung posisi baru dengan $x_n(time)=x_n(time-\Delta t) + v_n(time) \Delta t$.
	\State \hspace{0.5cm} \textbf{If}( $\Delta v \leq 10^{-5} \&\& a_n(time)\leq 10^{-5})$ \textbf{then}  
    \State \hspace{0.5cm} \hspace{0.5cm} \text{OUTPUT }Cetak hasil data $a_n, v_n, x_n$.
	\State \hspace{0.5cm} \hspace{0.5cm} \textbf{return}.
	\State \hspace{0.5cm} \textbf{End If}.
	\EndWhile
	\State \text{End}
 \EndProcedure
 \end{algorithmic}
 \caption{Prosedur simulasi dinamika lalu lintas menggunakan FVDM.}\label{Algo:FVDM}
\end{algorithm}
%
%\include{Hasil-dan-Pembahasan}
%
%\include{Kesimpulan}
%
\cleardoublepage
\addcontentsline{toc}{chapter}{Daftar Pustaka}
\bibliographystyle{acm} %harvard style
\bibliography{References}
%
%\pagebreak
\cleardoublepage
\addcontentsline{toc}{chapter}{Lampiran}
\chapter*{Lampiran}

\end{document}
